\documentclass[a4paper,11pt]{article}
        \usepackage[top=15mm,bottom=18mm,left=16mm,right=16mm]{geometry}
        \usepackage{fontspec}
        \usepackage{xeCJK}
        \setmainfont{Times New Roman}
        \setCJKmainfont{Yu Gothic}
        \setmonofont{Consolas}
        \setCJKmonofont{Yu Gothic}
        \usepackage{booktabs}
        \usepackage{longtable}
        \usepackage{tabularx}
        \usepackage{array}
        \usepackage{enumitem}
        \usepackage{hyperref}
        \usepackage{listings}
        \usepackage{xcolor}
        \hypersetup{
          colorlinks=true,
          linkcolor=blue,
          urlcolor=blue,
          pdftitle={ROS share / 相対パス解決},
          pdfauthor={Codex}
        }
        \lstset{
          basicstyle=\ttfamily\small,
          breaklines=true,
          columns=fullflexible,
          keepspaces=true,
          frame=single,
          framerule=0.2pt,
          rulecolor=\color{black},
          showstringspaces=false
        }
        \setlength{\parskip}{0.42em}
        \setlength{\parindent}{1em}
        \renewcommand{\arraystretch}{1.15}
        \title{09. ROS share / 相対パス解決}
        \author{solar\_ws0129-main}
        \date{2026-07-29}
        \begin{document}
        \maketitle

        \section*{対象}
        \begin{tabularx}{\linewidth}{>{\raggedright\arraybackslash}p{0.18\linewidth} X}
        \toprule
        項目 & 内容 \\
        \midrule
        ファイル & \path|mpc_solarcar/path_utils.py| \\
        種別 & Python \\
        区分 & config \\
        役割 & CWD、package share、repo root をまたいで path を実在ファイルへ解決する小さな基盤。 \\
        起動文脈 & Node 実行時の path ぶれを吸収する補助モジュール。 \\
        \bottomrule
        \end{tabularx}

        \section*{このファイルを読む理由}
        CWD、package share、repo root をまたいで path を実在ファイルへ解決する小さな基盤。

        \begin{itemize}[leftmargin=1.6em]
        \item ament の package share があればそちらを優先する。
\item インストール後の launch/node 実行でも同じ relative path を使えるようにする。
        \end{itemize}

        \section*{呼び出し関係}
        \subsection*{どこから呼ばれるか}
        \begin{itemize}[leftmargin=1.6em]
        \item \path|mpc_node.py|
\item \path|gps_sim_node.py|
\item \path|solar_state_node.py|
        \end{itemize}

        \subsection*{次に読むと理解が進むファイル}
        \begin{itemize}[leftmargin=1.6em]
        \item \path|mpc_solarcar/solar_profile.py|
        \end{itemize}

        \subsection*{同一リポジトリ内の主要依存}
        \begin{itemize}[leftmargin=1.6em]
        \item 同一リポジトリ内の明示 import は少ない。
        \end{itemize}

        \section*{ソース冒頭の把握}
        モジュール docstring または先頭意図の要約:

        モジュール docstring は明示されていない。

        \subsection*{代表コード断片}
        \begin{lstlisting}[language=Python]
REPO_ROOT = Path(__file__).resolve().parents[1]


def resolve_path(path: str, default_subdir: str = '') -> str:
    """Resolve a path relative to CWD or package share.

    - If absolute, return as-is.
    - If exists relative to CWD, return it.
    - Otherwise, resolve under <pkg_share>/<default_subdir> (or directly under share).
    """
    if path is None:
        return path
    path = os.path.expanduser(str(path))
    if os.path.isabs(path):
        return path
    if os.path.exists(path):
        return path
    if get_package_share_directory is not None:
        pkg_share = get_package_share_directory(PKG_NAME)
    else:
        pkg_share = os.fspath(REPO_ROOT)
    if default_subdir:
\end{lstlisting}


        \section*{主要構造の読み取り}
        主要関数は resolve\_path。

        \subsection*{主要クラス・関数}
\begin{longtable}{p{0.07\linewidth} p{0.16\linewidth} p{0.25\linewidth} p{0.40\linewidth}}
\toprule
行 & 種別 & 名前 & 補足 \\
\midrule
14 & function & \path|resolve_path| & args: path, default\_subdir \\
\bottomrule
\end{longtable}

        \subsection*{ファイルを上から読んだときの定義順}
            \begin{longtable}{p{0.07\linewidth} p{0.84\linewidth}}
            \toprule
            行 & 内容 \\
            \midrule
            4 & 例外処理を伴う try ブロックを実行する。 \\
10 & PKG\_NAME に 'mpc\_solarcar' の結果を代入する。 \\
11 & REPO\_ROOT に Path(\_\_file\_\_).resolve().parents[1] の結果を代入する。 \\
14 & 関数 resolve\_path を定義する。 \\
            \bottomrule
            \end{longtable}

        \subsection*{import 群の詳細}
            \begin{longtable}{p{0.07\linewidth} p{0.33\linewidth} p{0.50\linewidth}}
            \toprule
            行 & import 文 & このファイルで読む理由 \\
            \midrule
            1 & \path|import os| & パス、環境変数、プロセス外部状態を扱うため。 このファイル内での主な使用位置は L23, L24, L26, L31, L34, L35, L36, L37。 \\
2 & \path|from pathlib import Path| & ファイルやディレクトリを安全に扱うため。 このファイル内での主な使用位置は L11。 \\
5 & \path|from ament_index_python.packages import get_package_share_directory| & ament\_index\_python.packages から get\_package\_share\_directory を読み込み、このファイルの処理を組み立てるため。 このファイル内での主な使用位置は L7, L28, L29。 \\
            \bottomrule
            \end{longtable}

        \section*{関数・クラスを上から順に解説}
        \subsection*{L14 関数 \path|resolve_path|}
            \begin{itemize}[leftmargin=1.6em]
              \item 定義: \path|resolve_path(path, default_subdir)|
              \item docstring: Resolve a path relative to CWD or package share.

- If absolute, return as-is.
- If exists relative to CWD, return it.
- Otherwise, resolve under <pkg\_share>/<default\_subdir> (or directly under share).
              \item このブロックが直接呼ぶ主な関数・メソッド: exists, expanduser, fspath, get\_package\_share\_directory, isabs, join, startswith, str, strip
              \item 戻り値の要点: os.path.join(pkg\_share, path) / path / path / path
            \end{itemize}
            \begin{enumerate}[leftmargin=1.8em]
            \item 条件 path is None を判定し、真なら内部処理を行う。
\item path に os.path.expanduser(str(path)) の結果を代入する。
\item 条件 os.path.isabs(path) を判定し、真なら内部処理を行う。
\item 条件 os.path.exists(path) を判定し、真なら内部処理を行う。
\item 条件 get\_package\_share\_directory is not None を判定し、真なら内部処理を行う。
\item 条件 default\_subdir を判定し、真なら内部処理を行う。
\item os.path.join(pkg\_share, path) を返す。
            \end{enumerate}











        \section*{処理の流れ}
        \begin{enumerate}[leftmargin=1.7em]
        \item ソース中の主要関数を通じて処理を進める。
        \end{enumerate}

        \section*{このファイルの位置づけ}
        この資料群では、\path|mpc_solarcar/path_utils.py| を
        \textbf{config}の中核として扱う。
        したがって、ここで扱う処理は単独で完結せず、
        前段の profile / launch / telemetry と後段の planner / logger / report へ繋がる。

        \end{document}
